\section{Literature Review}

The literature relevant to this manuscript is narrower than the broader disaster-resilience field. The unresolved issue is not simply whether climate stress, institutions, or adaptation matter for rural livelihoods. The specific gap is whether international cocoa shocks are transmitted into producer-linked prices under macroeconomic conditions, whether that adjustment is distributed unevenly across the supply chain, and how local weather should be incorporated without being overstated as the dominant short-run price mechanism. The retained evidence is cocoa-led, with coffee studies used only where the manuscript explicitly invokes broader comparisons across tropical perennial cash crops cultivated under similar conditions. In that sense, the literature review extends the Introduction's argument rather than repeating it.

Vulnerability and food-system resilience concepts supply the framing for that gap. In these traditions, harm is not defined by exposure alone, but by exposure together with the sensitivity and adaptive capacity of the affected system \citep{Turner2003,Adger2006}. Resilience claims must also specify the relevant shock, actor, and timescale \citep{Tendall2015,Zurek2022}. These concepts justify treating cocoa price volatility as a resilience problem only insofar as it becomes producer exposure under constrained local conditions.

Evidence from cocoa-market integration shows that tighter integration with world markets can increase producer exposure to international volatility, while exchange-rate conditions shape cocoa price risk through distinct financial channels \citep{TsowouGayi2019,Jumah2001}. Broader price-transmission research further shows that pass-through along pricing chains is often partial and method-sensitive \citep{MeyerCramon2004,acosta_ihle_voncramon_2019}, with producer and consumer prices not necessarily adjusting in the same way \citep{JimenezRodriguez2022}. These studies justify modelling world cocoa prices and the COP/USD rate separately in the present manuscript. They do not, however, follow those shocks through producer-linked and downstream price layers in one empirical system.

The cocoa value-chain literature identifies a second missing piece. Producer concerns include price volatility, weak farmer organizations, dependence on few buyers, and incomplete alignment between interventions and producer priorities \citep{Mithofer2017}. More general value-chain and food-system research shows that governance structures and market power help determine how shocks are coordinated and absorbed across actors \citep{Gereffi2005,Merkle2021}. Comparative work on coffee should not be read as direct cocoa evidence, but it reinforces the broader mechanism that adaptive capacity in similar tropical tree-crop systems depends on support and coordination across the chain \citep{Morales2022,Verburg2019}. This combined evidence supports the idea that chain position matters for exposure, but it stops short of measuring whether downstream prices adjust differently from producer-linked cocoa prices in one aligned monthly design.

Climate-oriented cocoa studies identify the third component of the gap. Cocoa vulnerability has been analysed through climatic suitability, physiological sensitivity, and adaptation options \citep{Schroth2016,Laderach2013,Lahive2019}. Farm-system studies further show that agroforestry can strengthen buffering capacity and adaptation under cocoa climate stress \citep{Jacobi2013,Niether2020,Vaast2014}. Comparative coffee studies are used only to support broader cross-crop statements about spatial climate risk and adaptation planning \citep{Koh2020,Laderach2016}. Prior local cocoa modelling additionally supports the NASA POWER weather architecture used in the Methods section \citep{talero_sarmiento_etal_2025}. Within the vulnerability framework, these papers justify using localized weather indicators as contextual stress, yet they do not show that weather should replace benchmark transmission as the core short-run adjustment mechanism.

The manuscript therefore fills a specific gap left open across these strands. It combines international cocoa benchmarks, Colombian producer-linked prices, a European downstream chocolate price index, exchange-rate and oil controls, and localized weather indicators in one transparent empirical framework. The analytical gain comes from keeping those layers distinct: prices and macro controls identify transmission, downstream prices reveal how adjustment is distributed across the chain, and weather clarifies the local stress context in which transmitted shocks are experienced.
